\section{The exact point column at a simple VGIT wall}
\label{sec:simple-wall}

Let \(X_-\dashrightarrow X_+\) be a simple VGIT wall crossing satisfying the
hypotheses of \cite[Definition~3.2]{GuYuYu}: the chamber quotients and wall
quotient \(S\) are smooth and projective, and the normal weights at the wall
are \(\pm1\).  Assume \(r_+<r_-\) and write \(\nu=r_--r_+\).
\begin{proposition}[Pinned simple-wall input]
\label{prop:gyy-simple-wall-input}
Fix Gu--Yu--Yu, arXiv:2508.15770v1, and a simple wall as above. The following
are the only inputs from that source used in this section.
\begin{enumerate}[label=\textup{(G\arabic*)}]
\item Their decomposition theorem gives a pairing-preserving,
connection-intertwining decomposition over the exceptional Laurent completion
\cite[Theorem~6.2, p.~54]{GuYuYu}:
\begin{equation}\label{eq:simple-wall-qdm-decomposition}
 \Phi:\QDM(X_-)^{\mathrm{La,red}}
 \longrightarrow
 (\tau_+^{\mathrm{red}})^*\QDM(X_+)^{\mathrm{La,red}}
 \oplus\bigoplus_{j=0}^{\nu-1}
 (\zeta_j^{\mathrm{red}})^*\QDM(S)^{\mathrm{La,red}}.
\end{equation}
\item For common-open points \(p_\pm\), their Kirwan lift construction gives
a class \(a_p\) with \cite[Lemma~3.27, pp.~24--25]{GuYuYu}
\begin{equation}\label{eq:master-point-restrictions}
 \kappa_+(a_p)=[p_+],\qquad
 \kappa_-(a_p)=[p_-],\qquad
 a_p|_{F_0}=0.
\end{equation}
Here \(\kappa_\pm:H^*_{\mathbf C^*}(W)\to H^*(X_\pm)\) are the ordinary
cohomological Kirwan maps. The final equality is the point-support
specialization proved in that lemma: if
\([p_+]|_{\mathbf P(N_{F_0,+})}=f(h_+)\), then its proof identifies
\(a_p|_{F_0}=f(\lambda)\).  A common-open point misses the positive
exceptional locus, so \(f=0\).
\item At the extremal specialization,
\[
 \FT_{X_+}(\alpha)=\kappa_+(\alpha)
   +\sum_{k>0}f_{+,k}(\alpha)q_{\mathrm w}^k,
 \qquad f_{+,k}(\alpha)\in H^*(X_+)[z].
\]
On the negative side, if
\(\deg_\lambda(\alpha|_{F_0})\le r_{F_0,-}-1\), then
\(\FT_{X_-}(\alpha)=\kappa_-(\alpha)\) exactly at the same specialization.
The completed master source is finite free, and the negative Fourier map is
an isomorphism after this base change
\cite[Proposition~5.2, Lemmas~5.8 and~5.10, Proposition~5.9 and
Theorem~5.5, pp.~37, 42--44]{GuYuYu}.
\item Fourier transformation is covariant for the wall-frequency connection:
it carries multiplication by the equivariant wall class to the receiver's
wall logarithmic derivative together with quantum multiplication by its
Kirwan image and the receiver-coordinate derivative
\cite[Proposition~4.14(2), p.~32; Proposition~4.21, p.~36]{GuYuYu}.
\end{enumerate}
No sectorial or Gamma comparison is included in this input.
\end{proposition}

\begin{proof}
The ring in \textup{(G1)} is Laurent in a fractional wall variable and
complete in the remaining Novikov and bulk variables.  The four clauses are
the cited decomposition, Kirwan-lift, leading-term, and covariance statements,
respectively; no additional conclusion is drawn from them here.  We use only
the extremal specialization below.
\end{proof}

Let \(p_\pm\) be corresponding points in the common open set.  Gu--Yu--Yu's
lift \(a_p\) is the one in
Proposition~\ref{prop:gyy-simple-wall-input}\textup{(G2)}.

\begin{theorem}[Exact ambient point coordinate]
\label{thm:simple-wall-point-column}
After setting the master Novikov and bulk variables to zero while retaining
the wall variable, the ambient coordinate of
\eqref{eq:simple-wall-qdm-decomposition} satisfies
\begin{equation}\label{eq:simple-wall-point-column}
 \pr_{X_+}\Phi([p_-])=[p_+].
\end{equation}
This theorem makes no assertion about the centre coordinates of
\(\Phi([p_-])\).
\end{theorem}

\begin{proof}
Let \(q_{\mathrm w}=S_{F_0}\) denote the wall variable, let \(\lambda\) be the
equivariant wall frequency, and put \(D=\kappa_+(\lambda)\).
Proposition~\ref{prop:gyy-simple-wall-input}\textup{(G3)} applies to \(a_p\)
because its wall restriction is zero.  At the specialization under
consideration, it gives
\begin{equation}\label{eq:negative-fourier-point}
 \FT_{X_-}(a_p)=[p_-].
\end{equation}
Clause \textup{(G3)} gives only
\begin{equation}\label{eq:positive-fourier-point-tail}
 \FT_{X_+}(a_p)=[p_+]+O(q_{\mathrm w}).
\end{equation}

Apply \textup{(G3)} again to \(\lambda a_p\). Its wall restriction is zero and
\begin{equation}\label{eq:negative-frequency-kernel}
 \kappa_-(\lambda a_p)=\kappa_-(\lambda)[p_-]=0.
\end{equation}
by degree.  The finite-free/isomorphism clause of \textup{(G3)} survives base
change to the extremal specialization.
Equation~\eqref{eq:negative-frequency-kernel} and \textup{(G3)} therefore imply
that \(\lambda a_p\) is zero in the specialized completed source.  In
particular,
\begin{equation}\label{eq:positive-frequency-kernel}
 \FT_{X_+}(\lambda a_p)=0.
\end{equation}

The Fourier covariance of
Proposition~\ref{prop:gyy-simple-wall-input}\textup{(G4)} rewrites
\eqref{eq:positive-frequency-kernel} as
\begin{equation}\label{eq:positive-fourier-connection}
 \bigl(zq_{\mathrm w}\partial_{q_{\mathrm w}}+D\star_{\tau_+(q_{\mathrm w})}
 +(q_{\mathrm w}\partial_{q_{\mathrm w}}\tau_+(q_{\mathrm w}))
   \star_{\tau_+(q_{\mathrm w})}\bigr)
 \FT_{X_+}(a_p)=0.
\end{equation}
Here \(\tau_+(q_{\mathrm w})\) is the specialized receiver coordinate.  Write
\begin{equation}\label{eq:point-tail-expansions}
 \FT_{X_+}(a_p)=[p_+]+\sum_{k>0}q_{\mathrm w}^kv_k,
 \qquad
 \tau_+(q_{\mathrm w})=f(q_{\mathrm w})\mathbf1+\eta(q_{\mathrm w}),
\end{equation}
where \(f(0)=0=\eta(0)\) and \(\eta(q_{\mathrm w})\) has positive
cohomological degree.
Clause \textup{(G3)} gives \(v_k\in H^*(X_+)[z]\). Every nonconstant stable map in a
positive multiple of the wall fibre is contained in the exceptional locus
and misses \(p_+\).  Thus the positive-wall quantum corrections and
\(\eta(q_{\mathrm w})\) annihilate \([p_+]\).

We prove simultaneously that every \(v_k\) and every coefficient of \(f\)
vanishes.  Suppose this is known below order \(m>0\), and write
\([q_{\mathrm w}^m]f=c_m\).  The coefficient of \(q_{\mathrm w}^m\) in
\eqref{eq:positive-fourier-connection} is
\begin{equation}\label{eq:first-point-tail-coefficient}
 (zm\,\id+D\cup-)v_m+mc_m[p_+]=0.
\end{equation}
If \(v_m\ne0\), the highest power of \(z\) in \(zmv_m\) cannot be cancelled
by either \(D\cup v_m\) or the constant term \(mc_m[p_+]\), because \(v_m\)
is polynomial in \(z\).  Hence \(v_m=0\), and then \(c_m=0\).  Induction
kills the complete tail and the possible unit part of the receiver
coordinate, proving \eqref{eq:simple-wall-point-column}.
\end{proof}

\begin{warning}\label{warn:center-coordinates}
The continuous Fourier map applies
stationary phase to the transformed master fundamental solution.  Positive
master-Novikov terms can survive even though the classical restriction
\(a_p|_{F_0}\) is zero.  The argument proves the ambient coordinate exactly
because the wall-frequency derivative dies in the specialized completed
source; polynomiality from \textup{(G3)} then controls the possible receiver
unit correction.
\end{warning}

The following consequence isolates the extra analytic premise required to
turn \eqref{eq:simple-wall-point-column} into a Gamma-row statement.

\begin{hypothesis}[One-wall sectorial realization]
\label{hyp:one-wall-sectorial}
At a fixed nonzero wall parameter there is a pairing-preserving sectorial
realization of \eqref{eq:simple-wall-qdm-decomposition} on one common
admissible sector such that:
\begin{enumerate}[label=\textup{(\alph*)}]
\item the ambient formal summand realizes the intrinsic sectorial solution
  space of \(X_+\);
\item the wall summands realize the Gamma spans of the wall functors in the
  Shen--Shoemaker semiorthogonal decomposition;
\item the realized point row has ambient coordinate
  \eqref{eq:simple-wall-point-column}.
\end{enumerate}
\end{hypothesis}

\begin{corollary}[Conditional wall-local rank identity]
\label{cor:simple-wall-rank}
Under Hypothesis~\ref{hyp:one-wall-sectorial}, the point row restricts to the
point row of \(X_+\) on the ambient packet and vanishes on every wall packet.
In particular, for the whole generalized \(\zetaSix\) packet in this one
receiver,
\begin{equation}\label{eq:simple-wall-primary-boolean}
 b(X_-,\cP_6(X_-))=b(X_+,\cP_6(X_+)).
\end{equation}
\end{corollary}

\begin{proof}
The ambient assertion is Theorem~\ref{thm:simple-wall-point-column} together
with the pairing-preserving realization.  A wall functor has image supported
on the exceptional locus.  Pairing its Gamma class with the point class in
the common open gives zero Euler characteristic.  Shen--Shoemaker's oriented
Gamma asymptotics identify these supported classes with the wall blocks on
the chosen common sector \cite{ShenShoemaker}.  Hence the point row vanishes
there.  Restricting the resulting whole-block identity to the generalized
\(\zetaSix\) packet gives \eqref{eq:simple-wall-primary-boolean}.
\end{proof}

The hypothesis is not a reformulation of the cited formal theorems.
Gu--Yu--Yu work over an exceptional Laurent completion; Shen--Shoemaker's
asymptotic theorem is extremal and sectorial.  Identifying the intrinsic
large-radius point section with the full ambient packet across those frames
is additional data.
